%% ============================================================
%%  Chapter 3 — Database Design and Eloquent Models
%%  03_database_and_models.tex
%% ============================================================
\chapter{Database Design and Eloquent Models}
\label{chap:database}

\chapterintro{%
  This chapter documents the complete relational schema of the Clinic
  Management API. Every table is covered: its migration source code,
  column definitions, index strategy, and foreign key constraints are
  explained in full. Following the schema, every Eloquent model is
  presented with its relationship declarations. The chapter closes with
  the most comprehensive section: the Model Factory library, which is the
  foundation of the PHPUnit test suite. Special attention is given to the
  \phpclass{UserFactory} state pattern and the \phpclass{AppointmentFactory}
  static-counter technique for avoiding unique-key violations in batches.
}

%% ─────────────────────────────────────────────────────────────────────────
\section{Database Overview}
\label{sec:db-overview}
%% ─────────────────────────────────────────────────────────────────────────

The application uses a single MySQL 8.0 database with nine
application-owned tables. Six tables model the core business domain;
three infrastructure tables support framework concerns (OAuth tokens,
Laravel Cache, and the job queue).

\begin{table}[h!]
  \centering
  \caption{Complete table inventory.}
  \label{tab:table-inventory}
  \renewcommand{\arraystretch}{1.5}
  \begin{tabularx}{\linewidth}{@{}L{4cm} L{2.5cm} X@{}}
    \toprule
    \rowcolor{PrimaryBlue}
    \tblhdr{Table} & \tblhdr{Category} & \tblhdr{Purpose} \\
    \midrule
    \texttt{clinics}            & Domain   & One record per clinic branch \\
    \rowalt
    \texttt{users}              & Domain   & Staff accounts (Super Admin, Doctor, Assistant) \\
    \texttt{patients}           & Domain   & Patient registry, scoped per clinic \\
    \rowalt
    \texttt{doctor\_schedules}  & Domain   & Weekly availability windows per doctor \\
    \texttt{appointments}       & Domain   & Booked consultations \\
    \rowalt
    \texttt{prescriptions}      & Domain   & Consultation outcome records \\
    \texttt{prescription\_items}& Domain   & Individual medicine lines per prescription \\
    \rowalt
    \texttt{oauth\_access\_tokens} & Infrastructure & Laravel Passport token storage \\
    \texttt{cache}              & Infrastructure & Laravel Cache table driver \\
    \rowalt
    \texttt{jobs}               & Infrastructure & Laravel Queue table driver \\
    \bottomrule
  \end{tabularx}
\end{table}

%% ─────────────────────────────────────────────────────────────────────────
\section{The UUID Primary Key Strategy}
\label{sec:uuid-strategy}
%% ─────────────────────────────────────────────────────────────────────────

Every domain table uses a \texttt{CHAR(36)} UUID column as its primary
key rather than an auto-incrementing integer. This is enforced at the
Eloquent layer via the \texttt{HasUuids} trait applied to every model.

\subsection{How HasUuids Works}

When \texttt{HasUuids} is present on a model, Eloquent overrides the
key-generation strategy so that:

\begin{enumerate}
  \item The \texttt{incrementing} property is set to \texttt{false},
    preventing Eloquent from treating the primary key as a
    self-incrementing integer.
  \item The \texttt{keyType} property is set to \texttt{'string'} so
    that Eloquent performs string comparisons on the primary key rather
    than integer comparisons.
  \item Before any \texttt{INSERT}, the trait generates a new
    \texttt{Str::uuid()} value (RFC~4122 UUID v4) and assigns it to the
    primary key if one has not already been set.
\end{enumerate}

The migration column definition is simply:

\begin{lstlisting}[style=php, caption={UUID primary key column definition.}, label={lst:uuid-col}]
$table->uuid('id')->primary();
\end{lstlisting}

This expands to a \texttt{CHAR(36) NOT NULL PRIMARY KEY} column in MySQL.

\subsection{Why UUID Over Auto-Increment?}

\begin{table}[h!]
  \centering
  \caption{UUID vs.\ auto-increment integer trade-offs.}
  \label{tab:uuid-vs-int}
  \renewcommand{\arraystretch}{1.5}
  \begin{tabularx}{\linewidth}{@{}L{4cm} X X@{}}
    \toprule
    \rowcolor{PrimaryBlue}
    \tblhdr{Property} & \tblhdr{UUID} & \tblhdr{Auto-Increment} \\
    \midrule
    Exposable in API   & Safe --- reveals no row count or
                          insertion order
                       & Risky --- exposes row count and
                         sequential order \\
    \rowalt
    Distributed creation & Records can be created without a
                           central counter authority
                         & Requires a single write-authoritative
                           server \\
    Index performance  & Slightly slower B-tree insertions
                         due to random ordering
                       & Sequential insertions; optimal
                         B-tree performance \\
    \rowalt
    Join efficiency    & String comparison; slightly larger
                         than 4-byte integer
                       & Integer comparison; minimal storage \\
    Enumeration safety & Guessing valid IDs is computationally
                         infeasible ($2^{122}$ values)
                       & Sequential IDs are trivially enumerable \\
    \bottomrule
  \end{tabularx}
\end{table}

For a clinical API where resource IDs appear in URLs and API responses,
the security and distribution properties of UUID outweigh the minor
index-insertion performance cost.

%% ─────────────────────────────────────────────────────────────────────────
\section{PHP 8.1 Backed Enums}
\label{sec:enums}
%% ─────────────────────────────────────────────────────────────────────────

Three backed enums replace the magic strings that would otherwise be
scattered across the codebase. All three are used as Eloquent column
casts, meaning Eloquent automatically converts the raw database string
into the enum instance when a model is hydrated from a query result,
and converts the enum back to its string value on save.

\begin{lstlisting}[style=php, caption={app/Enums/UserRole.php.}, label={lst:enum-role}]
<?php

namespace App\Enums;

enum UserRole: string
{
    case SUPER_ADMIN = 'super_admin';
    case DOCTOR      = 'doctor';
    case ASSISTANT   = 'assistant';
}
\end{lstlisting}

\begin{lstlisting}[style=php, caption={app/Enums/AppointmentStatus.php.}, label={lst:enum-status}]
<?php

namespace App\Enums;

enum AppointmentStatus: string
{
    case PENDING   = 'pending';
    case CONFIRMED = 'confirmed';
    case COMPLETED = 'completed';
    case CANCELLED = 'cancelled';
}
\end{lstlisting}

\begin{lstlisting}[style=php, caption={app/Enums/Gender.php.}, label={lst:enum-gender}]
<?php

namespace App\Enums;

enum Gender: string
{
    case MALE   = 'male';
    case FEMALE = 'female';
}
\end{lstlisting}

\textbf{Why backed enums?} A plain (non-backed) PHP enum has no string
representation. A backed enum with \texttt{: string} associates each
case with an explicit string literal that is persisted to the database
and returned in API responses. This means that the migration column
definition (\texttt{\$table->enum(...)}) and the JSON API output can
both use human-readable lowercase strings (\texttt{"doctor"},
\texttt{"pending"}) without any manual serialisation code in the
controller or resource layers. Eloquent's enum casting handles the
round-trip conversion automatically.

\textbf{Type-safe comparisons.} Because the Eloquent attribute is
returned as a \phpclass{UserRole} instance rather than a bare string,
comparisons use enum case syntax:
\texttt{\$user->role === UserRole::DOCTOR} rather than
\texttt{\$user->role === 'doctor'}. This makes typos into compile-time
errors instead of silent logic bugs.

%% ─────────────────────────────────────────────────────────────────────────
\section{Entity Relationship Map}
\label{sec:erd}
%% ─────────────────────────────────────────────────────────────────────────

The following table describes every foreign key relationship in the
schema, its cardinality, and the \texttt{ON DELETE} behaviour enforced
at the database level.

\begin{longtable}{@{}L{3cm} L{3cm} C{1.5cm} L{3.5cm} X@{}}
  \toprule
  \rowcolor{PrimaryBlue}
  \tblhdr{From Table} & \tblhdr{Column} & \tblhdr{Card.}
    & \tblhdr{To Table} & \tblhdr{ON DELETE} \\
  \midrule
  \endfirsthead
  \toprule
  \rowcolor{PrimaryBlue}
  \tblhdr{From Table} & \tblhdr{Column} & \tblhdr{Card.}
    & \tblhdr{To Table} & \tblhdr{ON DELETE} \\
  \midrule
  \endhead
  \bottomrule
  \endfoot

  \texttt{users}
    & \texttt{clinic\_id}
    & $N$:1
    & \texttt{clinics}
    & SET NULL --- user record is preserved, clinic reference nulled \\
  \rowalt
  \texttt{patients}
    & \texttt{clinic\_id}
    & $N$:1
    & \texttt{clinics}
    & CASCADE --- deleting a clinic removes its patients \\
  \texttt{doctor\_schedules}
    & \texttt{doctor\_id}
    & $N$:1
    & \texttt{users}
    & CASCADE --- deleting a user removes their schedules \\
  \rowalt
  \texttt{doctor\_schedules}
    & \texttt{clinic\_id}
    & $N$:1
    & \texttt{clinics}
    & CASCADE \\
  \texttt{appointments}
    & \texttt{clinic\_id}
    & $N$:1
    & \texttt{clinics}
    & CASCADE \\
  \rowalt
  \texttt{appointments}
    & \texttt{doctor\_id}
    & $N$:1
    & \texttt{users}
    & CASCADE \\
  \texttt{appointments}
    & \texttt{patient\_id}
    & $N$:1
    & \texttt{patients}
    & CASCADE \\
  \rowalt
  \texttt{appointments}
    & \texttt{booked\_by}
    & $N$:1
    & \texttt{users}
    & CASCADE \\
  \texttt{prescriptions}
    & \texttt{appointment\_id}
    & 1:1
    & \texttt{appointments}
    & CASCADE \\
  \rowalt
  \texttt{prescriptions}
    & \texttt{doctor\_id}
    & $N$:1
    & \texttt{users}
    & CASCADE \\
  \texttt{prescriptions}
    & \texttt{patient\_id}
    & $N$:1
    & \texttt{patients}
    & CASCADE \\
  \rowalt
  \texttt{prescriptions}
    & \texttt{clinic\_id}
    & $N$:1
    & \texttt{clinics}
    & CASCADE \\
  \texttt{prescription\_items}
    & \texttt{prescription\_id}
    & $N$:1
    & \texttt{prescriptions}
    & CASCADE \\
\end{longtable}

\begin{notebox}[Cascade Semantics]
  The \texttt{CASCADE} strategy on appointment and prescription foreign
  keys means that deleting a \texttt{Clinic} record will recursively
  remove all its patients, appointments, prescriptions, and prescription
  items. This is appropriate for the clinic discontinuation use case
  (a clinic that closes has no need for its historical records to persist
  in isolation). However, the \texttt{SET NULL} strategy on
  \texttt{users.clinic\_id} deliberately preserves the user account,
  since a Super Administrator account has no clinic affiliation and
  doctor/assistant accounts may need to be reassigned to a replacement
  clinic.
\end{notebox}

%% ─────────────────────────────────────────────────────────────────────────
\section{Migrations: Core Domain Tables}
\label{sec:migrations}
%% ─────────────────────────────────────────────────────────────────────────

Laravel migrations are ordered by filename timestamp. The application
executes them in the following sequence: \texttt{clinics} first (because
\texttt{users} has a foreign key referencing it), then \texttt{users},
then all dependent tables. The \texttt{users} migration is split across
two files to resolve a chicken-and-egg dependency: the initial
\texttt{users} table is created without the foreign key constraint so
that the \texttt{clinics} table can be created first; a later migration
adds the constraint.

\subsection{clinics Table}

\begin{lstlisting}[
  style   = php,
  caption = {Migration: create\_clinics\_table (2026\_06\_05\_202852).},
  label   = {lst:mig-clinics},
]
Schema::create('clinics', function (Blueprint $table) {
    $table->uuid('id')->primary();
    $table->string('name');
    $table->string('address');
    $table->string('phone');
    $table->string('email');
    $table->boolean('is_active')->default(true);
    $table->timestamps();
});
\end{lstlisting}

\begin{table}[h!]
  \centering
  \caption{Column specification: \texttt{clinics}.}
  \label{tab:col-clinics}
  \renewcommand{\arraystretch}{1.5}
  \begin{tabularx}{\linewidth}{@{}L{3cm} L{3.2cm} L{1.8cm} X@{}}
    \toprule
    \rowcolor{PrimaryBlue}
    \tblhdr{Column} & \tblhdr{MySQL Type} & \tblhdr{Nullable} & \tblhdr{Notes} \\
    \midrule
    \texttt{id}        & \texttt{CHAR(36)} & No  & UUID PK, generated by \texttt{HasUuids} \\
    \rowalt
    \texttt{name}      & \texttt{VARCHAR(255)} & No  & Clinic trading name \\
    \texttt{address}   & \texttt{VARCHAR(255)} & No  & Physical street address \\
    \rowalt
    \texttt{phone}     & \texttt{VARCHAR(255)} & No  & Primary contact number \\
    \texttt{email}     & \texttt{VARCHAR(255)} & No  & Unique contact email \\
    \rowalt
    \texttt{is\_active}& \texttt{TINYINT(1)}   & No  & Default 1; soft-deactivation flag \\
    \texttt{created\_at} & \texttt{TIMESTAMP}  & Yes & Set by Eloquent on create \\
    \rowalt
    \texttt{updated\_at} & \texttt{TIMESTAMP}  & Yes & Set by Eloquent on update \\
    \bottomrule
  \end{tabularx}
\end{table}

\begin{notebox}
  The \texttt{clinics} table has no unique constraint on \texttt{name} or
  \texttt{email} at the migration level. Uniqueness for \texttt{email} is
  enforced in the \phpclass{StoreClinicRequest} and
  \phpclass{UpdateClinicRequest} Form Request validation rules.
  Adding a database-level unique index in a future migration is
  recommended for production deployments.
\end{notebox}

\subsection{users Table}

The \texttt{users} table is created in two stages. The initial migration
creates the table with \texttt{clinic\_id} as an \texttt{UNSIGNED BIGINT}
(the Laravel framework default) to avoid a forward-reference dependency.
A deferred migration later changes the column type to \texttt{CHAR(36)}
and adds the foreign key constraint.

\begin{lstlisting}[
  style   = php,
  caption = {Migration: create\_users\_table (0001\_01\_01\_000000) --- initial schema.},
  label   = {lst:mig-users-init},
]
Schema::create('users', function (Blueprint $table) {
    $table->uuid('id')->primary();
    $table->unsignedBigInteger('clinic_id')->nullable(); // FK added later
    $table->string('name');
    $table->string('email')->unique();
    $table->string('password');
    $table->string('phone')->nullable();
    $table->enum('role', ['super_admin', 'doctor', 'assistant']);
    $table->boolean('is_active')->default(true);
    $table->timestamps();
});
\end{lstlisting}

\begin{lstlisting}[
  style   = php,
  caption = {Migration: add\_foreign\_key\_to\_users\_clinic\_id (2026\_06\_11).},
  label   = {lst:mig-users-fk},
]
Schema::table('users', function (Blueprint $table) {
    // clinic_id was originally an unsignedBigInteger with no FK constraint.
    // Clinics use UUIDs (CHAR 36), so the column type must be changed
    // before the foreign key constraint can be added.
    $table->foreignUuid('clinic_id')->nullable()->change();
    $table->foreign('clinic_id')
          ->references('id')
          ->on('clinics')
          ->nullOnDelete();
});
\end{lstlisting}

\begin{warningbox}[Migration Ordering Dependency]
  The deferred FK migration \emph{must} run after both the
  \texttt{create\_clinics\_table} and the \texttt{create\_users\_table}
  migrations. The timestamp prefix \texttt{2026\_06\_11} ensures this
  ordering is respected. If the clinics table does not exist when this
  migration runs, MySQL will raise an error and the migration will fail.
  Always run \texttt{php artisan migrate} rather than executing
  individual migration files directly.
\end{warningbox}

\begin{table}[h!]
  \centering
  \caption{Column specification: \texttt{users}.}
  \label{tab:col-users}
  \renewcommand{\arraystretch}{1.5}
  \begin{tabularx}{\linewidth}{@{}L{3.2cm} L{3cm} L{1.8cm} X@{}}
    \toprule
    \rowcolor{PrimaryBlue}
    \tblhdr{Column} & \tblhdr{MySQL Type} & \tblhdr{Nullable} & \tblhdr{Notes} \\
    \midrule
    \texttt{id}         & \texttt{CHAR(36)}     & No  & UUID PK \\
    \rowalt
    \texttt{clinic\_id} & \texttt{CHAR(36)}     & Yes & FK $\rightarrow$ \texttt{clinics.id} ON DELETE SET NULL \\
    \texttt{name}       & \texttt{VARCHAR(255)} & No  & Full name \\
    \rowalt
    \texttt{email}      & \texttt{VARCHAR(255)} & No  & Unique; used as login credential \\
    \texttt{password}   & \texttt{VARCHAR(255)} & No  & Bcrypt hash; hidden from all serialisation \\
    \rowalt
    \texttt{phone}      & \texttt{VARCHAR(255)} & Yes & Optional contact number \\
    \texttt{role}       & \texttt{ENUM}         & No  & \texttt{'super\_admin'|'doctor'|'assistant'} \\
    \rowalt
    \texttt{is\_active} & \texttt{TINYINT(1)}   & No  & Default 1; deactivation blocks login \\
    \texttt{created\_at}& \texttt{TIMESTAMP}    & Yes & \\
    \rowalt
    \texttt{updated\_at}& \texttt{TIMESTAMP}    & Yes & \\
    \bottomrule
  \end{tabularx}
\end{table}

\begin{notebox}[No email\_verified\_at Column]
  The \texttt{users} table does not have the \texttt{email\_verified\_at}
  or \texttt{remember\_token} columns that Laravel includes in its default
  stub migration. These columns were deliberately omitted because the
  application does not implement email verification flows or web session
  (cookie-based) authentication. Both columns would be dead weight and
  their absence avoids a factory pitfall: the default Laravel
  \phpclass{UserFactory} sets \texttt{email\_verified\_at} to
  \texttt{now()}, which would cause a fatal \texttt{unknown column} error
  if copied without modification. The \phpclass{UserFactory} for this
  project does not set those fields.
\end{notebox}

\subsection{patients Table}

\begin{lstlisting}[
  style   = php,
  caption = {Migration: create\_patients\_table (2026\_06\_05\_202918).},
  label   = {lst:mig-patients},
]
Schema::create('patients', function (Blueprint $table) {
    $table->uuid('id')->primary();
    $table->foreignUuid('clinic_id')->constrained()->cascadeOnDelete();
    $table->string('name');
    $table->string('email')->nullable();
    $table->string('phone');
    $table->date('date_of_birth')->nullable();
    $table->enum('gender', ['male', 'female']);
    $table->text('address')->nullable();
    $table->timestamps();
});
\end{lstlisting}

The \texttt{patients} table is augmented by a separate FULLTEXT migration:

\begin{lstlisting}[
  style   = php,
  caption = {Migration: add\_fulltext\_index\_to\_patients (2026\_06\_11\_000002).},
  label   = {lst:mig-patients-ft},
]
Schema::table('patients', function (Blueprint $table) {
    $table->fullText(['name', 'phone']);
});
\end{lstlisting}

\textbf{Why a separate migration?} The FULLTEXT index migration is kept
separate for two reasons. First, it makes the index easy to drop and
recreate independently of the table definition (useful when
troubleshooting index performance). Second, it was added after the
initial schema was established, reflecting the real-world evolution of
the schema. A single migration can add both the table and the FULLTEXT
index, but splitting them documents the timeline of the design decision.

\textbf{FULLTEXT index on \texttt{(name, phone)}.} MySQL's FULLTEXT
engine indexes both columns as a composite virtual document. Searching
for a patient with
\texttt{->whereFullText(['name', 'phone'], 'ahmed 01001')} performs a
natural-language full-text search across both columns simultaneously.
This is several orders of magnitude faster than a
\texttt{LIKE '\%ahmed\%'} scan on large patient tables and requires
no additional query planner tuning.

\begin{importantbox}[MySQL Only]
  The FULLTEXT index is a MySQL-specific feature. The test suite is
  configured to use a real MySQL database (\texttt{DB\_CONNECTION=mysql})
  specifically to validate this functionality. Switching to SQLite for
  testing would silently make this index non-functional.
\end{importantbox}

\begin{table}[h!]
  \centering
  \caption{Column specification: \texttt{patients}.}
  \label{tab:col-patients}
  \renewcommand{\arraystretch}{1.5}
  \begin{tabularx}{\linewidth}{@{}L{3.2cm} L{3cm} L{1.8cm} X@{}}
    \toprule
    \rowcolor{PrimaryBlue}
    \tblhdr{Column} & \tblhdr{MySQL Type} & \tblhdr{Nullable} & \tblhdr{Notes} \\
    \midrule
    \texttt{id}            & \texttt{CHAR(36)}     & No  & UUID PK \\
    \rowalt
    \texttt{clinic\_id}    & \texttt{CHAR(36)}     & No  & FK $\rightarrow$ \texttt{clinics.id} CASCADE \\
    \texttt{name}          & \texttt{VARCHAR(255)} & No  & Indexed by FULLTEXT \\
    \rowalt
    \texttt{email}         & \texttt{VARCHAR(255)} & Yes & Optional \\
    \texttt{phone}         & \texttt{VARCHAR(255)} & No  & Indexed by FULLTEXT \\
    \rowalt
    \texttt{date\_of\_birth} & \texttt{DATE}       & Yes & Cast to \texttt{Carbon} by Eloquent \\
    \texttt{gender}        & \texttt{ENUM}         & No  & \texttt{'male'|'female'} \\
    \rowalt
    \texttt{address}       & \texttt{TEXT}         & Yes & \\
    \bottomrule
  \end{tabularx}
\end{table}

\subsection{doctor\_schedules Table}

\begin{lstlisting}[
  style   = php,
  caption = {Migration: create\_doctor\_schedules\_table (2026\_06\_05\_202954).},
  label   = {lst:mig-schedules},
]
Schema::create('doctor_schedules', function (Blueprint $table) {
    $table->uuid('id')->primary();
    $table->foreignUuid('doctor_id')->constrained('users')->cascadeOnDelete();
    $table->foreignUuid('clinic_id')->constrained()->cascadeOnDelete();
    $table->tinyInteger('day_of_week');  // 0 = Sunday ... 6 = Saturday
    $table->time('start_time');
    $table->time('end_time');
    $table->tinyInteger('slot_duration')->default(30); // minutes
    $table->boolean('is_active')->default(true);
    $table->timestamps();
});
\end{lstlisting}

\textbf{Column notes.}

\begin{itemize}
  \item \texttt{day\_of\_week} stores an integer from 0 to 6 following
    the \phpclass{Carbon} convention where 0 is Sunday and 6 is Saturday.
    This aligns with \texttt{Carbon::dayOfWeek} used in the slot
    computation algorithm in \phpclass{AppointmentService}.
  \item \texttt{start\_time} and \texttt{end\_time} use MySQL's native
    \texttt{TIME} type. MySQL returns these values as
    \texttt{HH:MM:SS} strings over the wire. As documented in
    Section~\ref{sec:time-normalisation}, the repository layer normalises
    these to \texttt{HH:MM} before any comparison with slot times.
  \item \texttt{slot\_duration} defaults to 30 minutes. The test suite
    uses 60-minute slots to keep slot boundary arithmetic simple.
  \item \texttt{is\_active} allows a schedule to be disabled temporarily
    without deleting it, supporting vacation or leave scenarios.
\end{itemize}

\begin{table}[h!]
  \centering
  \caption{Column specification: \texttt{doctor\_schedules}.}
  \label{tab:col-schedules}
  \renewcommand{\arraystretch}{1.5}
  \begin{tabularx}{\linewidth}{@{}L{3.2cm} L{3cm} L{1.8cm} X@{}}
    \toprule
    \rowcolor{PrimaryBlue}
    \tblhdr{Column} & \tblhdr{MySQL Type} & \tblhdr{Nullable} & \tblhdr{Notes} \\
    \midrule
    \texttt{id}             & \texttt{CHAR(36)}   & No & UUID PK \\
    \rowalt
    \texttt{doctor\_id}     & \texttt{CHAR(36)}   & No & FK $\rightarrow$ \texttt{users.id} CASCADE \\
    \texttt{clinic\_id}     & \texttt{CHAR(36)}   & No & FK $\rightarrow$ \texttt{clinics.id} CASCADE \\
    \rowalt
    \texttt{day\_of\_week}  & \texttt{TINYINT}    & No & 0=Sun, 1=Mon \ldots\ 6=Sat \\
    \texttt{start\_time}    & \texttt{TIME}       & No & Returns as \texttt{HH:MM:SS} from MySQL \\
    \rowalt
    \texttt{end\_time}      & \texttt{TIME}       & No & Returns as \texttt{HH:MM:SS} from MySQL \\
    \texttt{slot\_duration} & \texttt{TINYINT}    & No & Minutes per slot; default 30 \\
    \rowalt
    \texttt{is\_active}     & \texttt{TINYINT(1)} & No & Default 1 \\
    \bottomrule
  \end{tabularx}
\end{table}

\subsection{appointments Table}

The \texttt{appointments} table is the most constraint-rich table in
the schema, with four foreign keys, an enum column, and a composite
unique index.

\begin{lstlisting}[
  style   = php,
  caption = {Migration: create\_appointments\_table (2026\_06\_05\_203047).},
  label   = {lst:mig-appts},
]
Schema::create('appointments', function (Blueprint $table) {
    $table->uuid('id')->primary();
    $table->foreignUuid('clinic_id')->constrained()->cascadeOnDelete();
    $table->foreignUuid('doctor_id')->constrained('users')->cascadeOnDelete();
    $table->foreignUuid('patient_id')->constrained()->cascadeOnDelete();
    $table->foreignUuid('booked_by')->constrained('users')->cascadeOnDelete();

    $table->date('appointment_date');
    $table->time('start_time');
    $table->time('end_time');
    $table->enum('status', ['pending', 'confirmed', 'cancelled', 'completed'])
          ->default('pending');
    $table->text('notes')->nullable();
    $table->timestamps();

    // Last-resort double-booking guard at the database level.
    // The application-layer cache lock is the primary guard.
    $table->unique(['doctor_id', 'appointment_date', 'start_time']);
});
\end{lstlisting}

\textbf{The composite unique index} on
\texttt{(doctor\_id, appointment\_date, start\_time)} is the database's
final defence against double-booking. The application layer uses a
\texttt{Cache::lock()} (see Section~\ref{sec:service-layer}) as the
primary concurrency guard, but the database index provides a hard
guarantee in all scenarios, including direct database access outside
the application.

This index also has a significant implication for the test factory,
which is addressed in Section~\ref{sec:appt-factory}: creating two
\texttt{Appointment} records in the same test with the same doctor, date,
and start time will raise a duplicate key error unless the test provides
distinct start times.

\textbf{The \texttt{booked\_by} column} records the \texttt{users.id}
of the Assistant who created the appointment record. It is a foreign key
to \texttt{users} with a non-standard local column name, requiring the
explicit \texttt{constrained('users')} call to tell Laravel which table
to reference.

\begin{table}[h!]
  \centering
  \caption{Column specification: \texttt{appointments}.}
  \label{tab:col-appointments}
  \renewcommand{\arraystretch}{1.5}
  \begin{tabularx}{\linewidth}{@{}L{3.5cm} L{3cm} L{1.8cm} X@{}}
    \toprule
    \rowcolor{PrimaryBlue}
    \tblhdr{Column} & \tblhdr{MySQL Type} & \tblhdr{Nullable} & \tblhdr{Notes} \\
    \midrule
    \texttt{id}               & \texttt{CHAR(36)} & No  & UUID PK \\
    \rowalt
    \texttt{clinic\_id}       & \texttt{CHAR(36)} & No  & FK $\rightarrow$ \texttt{clinics.id} CASCADE \\
    \texttt{doctor\_id}       & \texttt{CHAR(36)} & No  & FK $\rightarrow$ \texttt{users.id} CASCADE \\
    \rowalt
    \texttt{patient\_id}      & \texttt{CHAR(36)} & No  & FK $\rightarrow$ \texttt{patients.id} CASCADE \\
    \texttt{booked\_by}       & \texttt{CHAR(36)} & No  & FK $\rightarrow$ \texttt{users.id} CASCADE \\
    \rowalt
    \texttt{appointment\_date}& \texttt{DATE}     & No  & Cast to \texttt{Carbon} by Eloquent \\
    \texttt{start\_time}      & \texttt{TIME}     & No  & Part of composite unique index \\
    \rowalt
    \texttt{end\_time}        & \texttt{TIME}     & No  & Computed by service; never from client \\
    \texttt{status}           & \texttt{ENUM}     & No  & Default \texttt{'pending'} \\
    \rowalt
    \texttt{notes}            & \texttt{TEXT}     & Yes & Optional booking notes \\
    \bottomrule
  \end{tabularx}
\end{table}

\subsection{prescriptions Table}

\begin{lstlisting}[
  style   = php,
  caption = {Migration: create\_prescriptions\_table (2026\_06\_05\_203107).},
  label   = {lst:mig-rx},
]
Schema::create('prescriptions', function (Blueprint $table) {
    $table->uuid('id')->primary();
    $table->foreignUuid('appointment_id')->constrained()->cascadeOnDelete();
    $table->foreignUuid('doctor_id')->constrained('users')->cascadeOnDelete();
    $table->foreignUuid('patient_id')->constrained()->cascadeOnDelete();
    $table->foreignUuid('clinic_id')->constrained()->cascadeOnDelete();
    $table->text('diagnosis');
    $table->text('notes')->nullable();
    $table->timestamps();
});
\end{lstlisting}

The prescription table denormalises \texttt{doctor\_id},
\texttt{patient\_id}, and \texttt{clinic\_id} even though all three
could be derived from the linked \texttt{appointment}. This is a
deliberate design choice: it allows efficient prescription queries by
doctor (\texttt{WHERE doctor\_id = ?}) and patient without joining
through the \texttt{appointments} table, and it provides an explicit
audit trail of who wrote what for whom, independent of any future
appointment archival strategy.

\subsection{prescription\_items Table}

\begin{lstlisting}[
  style   = php,
  caption = {Migration: create\_prescription\_items\_table (2026\_06\_05\_203128).},
  label   = {lst:mig-rx-items},
]
Schema::create('prescription_items', function (Blueprint $table) {
    $table->uuid('id')->primary();
    $table->foreignUuid('prescription_id')->constrained()->cascadeOnDelete();
    $table->string('medicine_name');
    $table->string('dosage');
    $table->string('frequency');
    $table->string('duration');
    $table->text('notes')->nullable();
    $table->timestamps();
});
\end{lstlisting}

All fields except \texttt{notes} are non-nullable. The
\phpclass{StorePrescriptionRequest} enforces this at validation time:
each item in the \texttt{items} array must provide
\texttt{medicine\_name}, \texttt{dosage}, \texttt{frequency}, and
\texttt{duration}. The \texttt{notes} field is optional per item,
distinct from the top-level prescription \texttt{notes}.

%% ─────────────────────────────────────────────────────────────────────────
\section{Eloquent Models}
\label{sec:models}
%% ─────────────────────────────────────────────────────────────────────────

\subsection{The User Model}

\begin{lstlisting}[
  style   = php,
  caption = {app/Models/User.php.},
  label   = {lst:model-user},
]
<?php

namespace App\Models;

use App\Enums\UserRole;
use Database\Factories\UserFactory;
use Illuminate\Database\Eloquent\Concerns\HasUuids;
use Illuminate\Database\Eloquent\Factories\HasFactory;
use Illuminate\Database\Eloquent\Relations\BelongsTo;
use Illuminate\Database\Eloquent\Relations\HasMany;
use Illuminate\Foundation\Auth\User as Authenticatable;
use Illuminate\Notifications\Notifiable;
use Laravel\Passport\HasApiTokens;

class User extends Authenticatable
{
    /** @use HasFactory<UserFactory> */
    use HasApiTokens, HasFactory, Notifiable, HasUuids;

    protected $fillable = [
        'clinic_id', 'name', 'email', 'password',
        'phone', 'role', 'is_active',
    ];

    protected $hidden = ['password', 'remember_token'];

    protected $casts = [
        'password'  => 'hashed',
        'is_active' => 'boolean',
        'role'      => UserRole::class,
    ];

    public function clinic(): BelongsTo
    {
        return $this->belongsTo(Clinic::class);
    }

    public function schedules(): HasMany
    {
        return $this->hasMany(DoctorSchedule::class);
    }
}
\end{lstlisting}

\textbf{Key observations on the User model.}

\begin{itemize}
  \item \textbf{\texttt{HasApiTokens}} is the Passport trait. It adds the
    \texttt{createToken()} method and the \texttt{token()} accessor used
    by \phpclass{AuthService} to issue and revoke access tokens.

  \item \textbf{\texttt{extends Authenticatable}} rather than plain
    \texttt{Model}. This inheritance chain (\texttt{User} $\to$
    \texttt{Authenticatable} $\to$ \texttt{Model}) satisfies Laravel's
    authentication guard contract without requiring manual interface
    implementation.

  \item \textbf{\texttt{'password' => 'hashed'}} is a Laravel 10+
    cast that automatically bcrypt-hashes any value assigned to the
    \texttt{password} attribute, removing the need for explicit
    \texttt{bcrypt()} or \texttt{Hash::make()} calls when setting a
    password via mass assignment.

  \item \textbf{\texttt{'role' => UserRole::class}} is the enum cast.
    When Eloquent loads a user from MySQL, the raw string
    \texttt{"doctor"} in the \texttt{role} column is automatically
    converted to the \texttt{UserRole::DOCTOR} enum case. On save, the
    enum case is converted back to its \texttt{->value} string.

  \item \textbf{\texttt{schedules(): HasMany}} is the only
    \texttt{HasMany} relationship on the User model. The model does not
    declare \texttt{hasMany(Appointment::class)} because appointments
    are always queried through the repository with explicit where clauses
    rather than through the Eloquent relationship chain, avoiding N+1
    query exposure.
\end{itemize}

\subsection{The Clinic Model}

\begin{lstlisting}[
  style   = php,
  caption = {app/Models/Clinic.php.},
  label   = {lst:model-clinic},
]
<?php

namespace App\Models;

use Illuminate\Database\Eloquent\Concerns\HasUuids;
use Illuminate\Database\Eloquent\Factories\HasFactory;
use Illuminate\Database\Eloquent\Model;
use Illuminate\Database\Eloquent\Relations\HasMany;

class Clinic extends Model
{
    use HasFactory, HasUuids;

    protected $fillable = [
        'name', 'address', 'phone', 'email', 'is_active',
    ];

    protected $casts = ['is_active' => 'boolean'];

    public function doctors(): HasMany
    {
        return $this->hasMany(User::class);
    }

    public function appointments(): HasMany
    {
        return $this->hasMany(Appointment::class);
    }

    public function patients(): HasMany
    {
        return $this->hasMany(Patient::class);
    }
}
\end{lstlisting}

The \texttt{Clinic} model is the aggregate root: every other domain
entity has a foreign key pointing back to a clinic. The three
\texttt{HasMany} relationships declared here are available for eager
loading in dashboard aggregate queries but are not used directly in
per-request resource retrieval, which goes through scoped repository
queries instead.

\subsection{The Patient Model}

\begin{lstlisting}[
  style   = php,
  caption = {app/Models/Patient.php.},
  label   = {lst:model-patient},
]
<?php

namespace App\Models;

use App\Enums\Gender;
use Illuminate\Database\Eloquent\Concerns\HasUuids;
use Illuminate\Database\Eloquent\Factories\HasFactory;
use Illuminate\Database\Eloquent\Model;
use Illuminate\Database\Eloquent\Relations\BelongsTo;
use Illuminate\Database\Eloquent\Relations\HasMany;

class Patient extends Model
{
    use HasFactory, HasUuids;

    protected $fillable = [
        'clinic_id', 'name', 'email', 'phone',
        'date_of_birth', 'gender', 'address',
    ];

    protected $casts = [
        'date_of_birth' => 'date',
        'gender'        => Gender::class,
    ];

    public function clinic(): BelongsTo
    {
        return $this->belongsTo(Clinic::class);
    }

    public function appointments(): HasMany
    {
        return $this->hasMany(Appointment::class);
    }
}
\end{lstlisting}

The \texttt{date} cast on \texttt{date\_of\_birth} converts the MySQL
\texttt{DATE} string into a \texttt{Carbon} instance on hydration,
enabling age calculation via \texttt{\$patient->date\_of\_birth->age}
and date comparison without manual string parsing.

\subsection{The Appointment Model}

\begin{lstlisting}[
  style   = php,
  caption = {app/Models/Appointment.php.},
  label   = {lst:model-appt},
]
<?php

namespace App\Models;

use App\Enums\AppointmentStatus;
use Illuminate\Database\Eloquent\Concerns\HasUuids;
use Illuminate\Database\Eloquent\Factories\HasFactory;
use Illuminate\Database\Eloquent\Model;
use Illuminate\Database\Eloquent\Relations\BelongsTo;
use Illuminate\Database\Eloquent\Relations\HasOne;

class Appointment extends Model
{
    use HasFactory, HasUuids;

    protected $fillable = [
        'clinic_id', 'doctor_id', 'patient_id', 'booked_by',
        'appointment_date', 'start_time', 'end_time',
        'status', 'notes',
    ];

    protected $attributes = ['status' => 'pending'];

    protected $casts = [
        'appointment_date' => 'date',
        'status'           => AppointmentStatus::class,
    ];

    public function clinic(): BelongsTo
    {
        return $this->belongsTo(Clinic::class);
    }

    public function doctor(): BelongsTo
    {
        return $this->belongsTo(User::class, 'doctor_id');
    }

    public function patient(): BelongsTo
    {
        return $this->belongsTo(Patient::class);
    }

    public function bookedBy(): BelongsTo
    {
        return $this->belongsTo(User::class, 'booked_by');
    }

    public function prescription(): HasOne
    {
        return $this->hasOne(Prescription::class);
    }
}
\end{lstlisting}

\textbf{Key observations.}

\begin{itemize}
  \item \textbf{\texttt{\$attributes = ['status' => 'pending']}} sets
    the default value at the Eloquent level. Without this, a newly
    instantiated \texttt{Appointment} model would have a null status
    until it is saved, which could cause a null pointer dereference if
    status is read before the first database round-trip.

  \item \textbf{\texttt{doctor()} and \texttt{bookedBy()} both point to
    \texttt{User::class}} but use different foreign keys. Eloquent
    cannot infer the correct foreign key column from the method name
    alone when it differs from the \texttt{user\_id} convention, so
    the second argument (\texttt{'doctor\_id'} and
    \texttt{'booked\_by'}) must be explicit.

  \item \textbf{\texttt{prescription(): HasOne}} is a one-to-one
    relationship. An appointment has at most one prescription. The
    \phpclass{PrescriptionController} enforces this at the business
    logic level; the database does not have a unique constraint on
    \texttt{prescriptions.appointment\_id}, making this a future
    hardening opportunity.
\end{itemize}

\subsection{The DoctorSchedule Model}

\begin{lstlisting}[
  style   = php,
  caption = {app/Models/DoctorSchedule.php.},
  label   = {lst:model-schedule},
]
<?php

namespace App\Models;

use Illuminate\Database\Eloquent\Concerns\HasUuids;
use Illuminate\Database\Eloquent\Factories\HasFactory;
use Illuminate\Database\Eloquent\Model;
use Illuminate\Database\Eloquent\Relations\BelongsTo;

class DoctorSchedule extends Model
{
    use HasFactory, HasUuids;

    protected $fillable = [
        'doctor_id', 'clinic_id', 'day_of_week',
        'start_time', 'end_time', 'slot_duration', 'is_active',
    ];

    protected $casts = ['is_active' => 'boolean'];

    public function doctor(): BelongsTo
    {
        return $this->belongsTo(User::class, 'doctor_id');
    }

    public function clinic(): BelongsTo
    {
        return $this->belongsTo(Clinic::class);
    }
}
\end{lstlisting}

\subsection{The Prescription and PrescriptionItem Models}

\begin{lstlisting}[
  style   = php,
  caption = {app/Models/Prescription.php.},
  label   = {lst:model-rx},
]
<?php

namespace App\Models;

use Illuminate\Database\Eloquent\Concerns\HasUuids;
use Illuminate\Database\Eloquent\Factories\HasFactory;
use Illuminate\Database\Eloquent\Model;
use Illuminate\Database\Eloquent\Relations\BelongsTo;
use Illuminate\Database\Eloquent\Relations\HasMany;

class Prescription extends Model
{
    use HasFactory, HasUuids;

    protected $fillable = [
        'appointment_id', 'doctor_id', 'patient_id',
        'clinic_id', 'diagnosis', 'notes',
    ];

    public function appointment(): BelongsTo
    {
        return $this->belongsTo(Appointment::class);
    }

    public function doctor(): BelongsTo
    {
        return $this->belongsTo(User::class, 'doctor_id');
    }

    public function patient(): BelongsTo
    {
        return $this->belongsTo(Patient::class);
    }

    public function clinic(): BelongsTo
    {
        return $this->belongsTo(Clinic::class);
    }

    public function items(): HasMany
    {
        return $this->hasMany(PrescriptionItem::class);
    }
}
\end{lstlisting}

\begin{lstlisting}[
  style   = php,
  caption = {app/Models/PrescriptionItem.php.},
  label   = {lst:model-rx-item},
]
<?php

namespace App\Models;

use Illuminate\Database\Eloquent\Concerns\HasUuids;
use Illuminate\Database\Eloquent\Factories\HasFactory;
use Illuminate\Database\Eloquent\Model;
use Illuminate\Database\Eloquent\Relations\BelongsTo;

class PrescriptionItem extends Model
{
    use HasFactory, HasUuids;

    protected $fillable = [
        'prescription_id', 'medicine_name', 'dosage',
        'frequency', 'duration', 'notes',
    ];

    public function prescription(): BelongsTo
    {
        return $this->belongsTo(Prescription::class);
    }
}
\end{lstlisting}

%% ─────────────────────────────────────────────────────────────────────────
\section{Model Factories}
\label{sec:factories}
%% ─────────────────────────────────────────────────────────────────────────

Model Factories are the cornerstone of the PHPUnit test suite. They are
PHP classes that extend
\texttt{Illuminate\textbackslash{}Database\textbackslash{}Eloquent\textbackslash{}Factories\textbackslash{}Factory}
and provide a \texttt{definition()} method that returns an array of
default attribute values for one model instance. The Laravel testing
framework provides \texttt{Factory::create()} (persists to DB) and
\texttt{Factory::make()} (in-memory only) for instantiating models from
these definitions.

The factory system supports three essential testing capabilities:

\begin{enumerate}
  \item \textbf{Realistic random data.} The \texttt{fake()} helper
    (backed by the Faker library) generates realistic names, emails,
    phone numbers, and dates, avoiding hardcoded strings that accumulate
    unique-constraint collisions across test runs.
  \item \textbf{Automatic relationship resolution.} A factory attribute
    can reference another factory (e.g.\ \texttt{'clinic\_id' =>
    Clinic::factory()}), which Laravel automatically resolves to a
    created record's UUID when the factory runs.
  \item \textbf{State methods.} Fluent state methods allow tests to
    express intent clearly and compose multiple states without test
    helper duplication (e.g.\
    \texttt{User::factory()->doctor()->forClinic(\$clinic)->inactive()->create()}).
\end{enumerate}

\subsection{UserFactory --- The State Pattern in Detail}
\label{sec:user-factory}

\phpclass{UserFactory} is the most important factory in the codebase
because \texttt{User} records are required by almost every test. The
factory must support all three role types, clinic affiliation, and the
active/inactive toggle.

\begin{lstlisting}[
  style   = php,
  caption = {database/factories/UserFactory.php --- complete source.},
  label   = {lst:factory-user},
]
<?php

namespace Database\Factories;

use App\Enums\UserRole;
use App\Models\Clinic;
use App\Models\User;
use Illuminate\Database\Eloquent\Factories\Factory;
use Illuminate\Support\Facades\Hash;

/**
 * @extends Factory<User>
 */
class UserFactory extends Factory
{
    // Static property: the hashed password is computed once and reused
    // across all factory calls within a single test run. This avoids
    // bcrypt's intentional slowness (typically 100ms per hash) from
    // multiplying across dozens of user creation calls.
    protected static ?string $password;

    public function definition(): array
    {
        return [
            'clinic_id' => null,
            'name'      => fake()->name(),
            'email'     => fake()->unique()->safeEmail(),
            'password'  => static::$password ??= Hash::make('password'),
            'phone'     => fake()->phoneNumber(),
            'role'      => UserRole::DOCTOR,       // default role
            'is_active' => true,
        ];
    }

    /** State: Super Administrator — global scope, no clinic binding */
    public function superAdmin(): static
    {
        return $this->state([
            'role'      => UserRole::SUPER_ADMIN,
            'clinic_id' => null,
        ]);
    }

    /** State: Doctor — default role, requires clinic via forClinic() */
    public function doctor(): static
    {
        return $this->state(['role' => UserRole::DOCTOR]);
    }

    /** State: Assistant — front-desk staff, requires clinic via forClinic() */
    public function assistant(): static
    {
        return $this->state(['role' => UserRole::ASSISTANT]);
    }

    /** State: Bind the user to a specific Clinic instance */
    public function forClinic(Clinic $clinic): static
    {
        return $this->state(['clinic_id' => $clinic->id]);
    }

    /** State: Deactivated account — blocks login via AccountDeactivatedException */
    public function inactive(): static
    {
        return $this->state(['is_active' => false]);
    }
}
\end{lstlisting}

\subsubsection{The \texttt{static::\$password} Optimisation}

The single most impactful performance technique in the factory is the
static property cache on the password hash:

\begin{lstlisting}[style=php, caption={Password hash memoisation.}, label={lst:pw-cache}]
protected static ?string $password;

// Inside definition():
'password' => static::$password ??= Hash::make('password'),
\end{lstlisting}

\textbf{Why this matters.} \texttt{Hash::make()} uses the bcrypt
algorithm with a cost factor of 12 (the Laravel default). At this cost,
a single bcrypt operation takes approximately 80--120 milliseconds on
typical hardware. The test suite creates 73+ user records across all
test cases. Without caching, bcrypt would run 73+ times, adding up to
\emph{6--9 seconds} of pure hashing time to the test run.

\textbf{How it works.} The \texttt{??=} null-coalescing assignment
operator checks whether \texttt{static::\$password} is null. On the
first call within a PHP process, it is null, so the right-hand side
(\texttt{Hash::make('password')}) is evaluated and the result is stored
in \texttt{static::\$password}. On every subsequent call, the
stored hash is returned immediately without re-hashing.

\textbf{The \texttt{static::} vs \texttt{self::} distinction.} Using
\texttt{static::\$password} instead of \texttt{self::\$password} ensures
that if a subclass of \texttt{UserFactory} is created in the future, it
uses its own static property rather than sharing the parent class's cache.
This is the Late Static Binding (LSB) semantics of PHP 5.3+.

\textbf{Test isolation.} Because the cache is a \emph{static} property
(class-level, not instance-level), it persists across factory instances
but is reset between PHPUnit processes. Within a single test run, all
test cases share the same hash. This is safe because the password value
used by factories (\texttt{'password'}) is a constant known to the
test suite and is never used in production.

\subsubsection{State Method Mechanics}

Each state method calls \texttt{\$this->state(array \$attributes)}. The
\texttt{state()} method returns a new \texttt{Factory} instance (the
factory is immutable) with the given attributes merged over the defaults
from \texttt{definition()}. States are applied in order and later states
override earlier ones.

This means:

\begin{lstlisting}[style=php, caption={State composition example.}, label={lst:state-compose}]
// All three states are applied in sequence.
// forClinic() overrides the clinic_id that superAdmin() set to null.
// inactive() overrides the is_active that definition() set to true.
$user = User::factory()
    ->superAdmin()
    ->forClinic($clinic) // overrides clinic_id => null from superAdmin()
    ->inactive()         // overrides is_active => true from definition()
    ->create();
// Result: role=SUPER_ADMIN, clinic_id=$clinic->id, is_active=false
\end{lstlisting}

The composability of state methods is what allows tests to express
complex fixture scenarios in a single readable line without
intermediate variable assignments or boilerplate arrays.

\subsubsection{The \texttt{forClinic()} State and Type Safety}

\begin{lstlisting}[style=php, caption={forClinic() state.}, label={lst:for-clinic}]
public function forClinic(Clinic $clinic): static
{
    return $this->state(['clinic_id' => $clinic->id]);
}
\end{lstlisting}

The \texttt{forClinic()} method accepts a \texttt{Clinic} \emph{model
instance}, not a raw string UUID. This type declaration has two benefits:

\begin{enumerate}
  \item \textbf{IDE autocomplete.} Passing a \texttt{Clinic} instance
    to \texttt{forClinic()} allows the IDE to type-check the argument
    and flag any mismatched types at development time, not at runtime.
  \item \textbf{Explicit UUID extraction.} \texttt{\$clinic->id} is the
    UUID string. Accepting the model instance rather than the string
    prevents callers from accidentally passing an integer ID or a
    differently typed identifier.
\end{enumerate}

In tests, this is always used in combination with the clinic created in
\texttt{setUp()}:

\begin{lstlisting}[style=php, caption={Typical setUp() fixture pattern.}, label={lst:setup-pattern}]
protected function setUp(): void
{
    parent::setUp();
    $this->clinic    = Clinic::factory()->create();
    $this->doctor    = User::factory()->doctor()->forClinic($this->clinic)->create();
    $this->assistant = User::factory()->assistant()->forClinic($this->clinic)->create();
    $this->actingAsPassport($this->assistant);
}
\end{lstlisting}

\subsection{ClinicFactory}

\begin{lstlisting}[
  style   = php,
  caption = {database/factories/ClinicFactory.php.},
  label   = {lst:factory-clinic},
]
<?php

namespace Database\Factories;

use App\Models\Clinic;
use Illuminate\Database\Eloquent\Factories\Factory;

/** @extends Factory<Clinic> */
class ClinicFactory extends Factory
{
    protected $model = Clinic::class;

    public function definition(): array
    {
        return [
            'name'      => fake()->company(),
            'address'   => fake()->streetAddress(),
            'phone'     => fake()->numerify('01#########'),
            'email'     => fake()->unique()->companyEmail(),
            'is_active' => true,
        ];
    }

    /** State: inactive clinic (soft-discontinued branch) */
    public function inactive(): static
    {
        return $this->state(['is_active' => false]);
    }
}
\end{lstlisting}

\textbf{Notes on \texttt{ClinicFactory}.}

\begin{itemize}
  \item \texttt{fake()->numerify('01\#\#\#\#\#\#\#\#\#\#')} generates
    an 11-digit phone number with \texttt{01} prefix, matching the
    Egyptian mobile number format used in the domain. The \texttt{\#}
    placeholder is replaced with a random digit.
  \item \texttt{fake()->unique()->companyEmail()} ensures no two
    test clinics share an email address within a test run. The
    \texttt{unique()} modifier maintains a per-faker-instance registry
    of previously generated values and retries on collision.
  \item \textbf{The \texttt{protected \$model} declaration} is required
    on concrete factory subclasses for models whose class name cannot
    be automatically inferred from the factory class name using
    Laravel's naming convention. Since \phpclass{ClinicFactory}
    correctly maps to \phpclass{Clinic} by convention, this is
    technically redundant but is included for clarity.
\end{itemize}

\subsection{PatientFactory}

\begin{lstlisting}[
  style   = php,
  caption = {database/factories/PatientFactory.php.},
  label   = {lst:factory-patient},
]
<?php

namespace Database\Factories;

use App\Enums\Gender;
use App\Models\Clinic;
use App\Models\Patient;
use Illuminate\Database\Eloquent\Factories\Factory;

/** @extends Factory<Patient> */
class PatientFactory extends Factory
{
    protected $model = Patient::class;

    public function definition(): array
    {
        return [
            'clinic_id'     => Clinic::factory(),
            'name'          => fake()->name(),
            'email'         => fake()->unique()->safeEmail(),
            'phone'         => fake()->numerify('01#########'),
            'date_of_birth' => fake()->date('Y-m-d', '-18 years'),
            'gender'        => fake()->randomElement([Gender::MALE, Gender::FEMALE]),
            'address'       => fake()->address(),
        ];
    }
}
\end{lstlisting}

\textbf{Notes on \texttt{PatientFactory}.}

\begin{itemize}
  \item \textbf{\texttt{'clinic\_id' => Clinic::factory()}} demonstrates
    factory chaining. When \texttt{PatientFactory::create()} is called
    without an explicit \texttt{clinic\_id}, Laravel automatically
    calls \texttt{Clinic::factory()->create()} to produce a Clinic
    record and uses its UUID as the \texttt{clinic\_id}. When a test
    provides \texttt{clinic\_id} explicitly (e.g.\
    \texttt{Patient::factory()->create(['clinic\_id' => \$this->clinic->id])}),
    the chained factory is skipped entirely.

  \item \textbf{\texttt{fake()->date('Y-m-d', '-18 years')}} generates
    a birth date for a person who is at least 18 years old (the
    \texttt{'-18 years'} argument is the maximum date, not minimum
    age). This ensures test patients are always legal adults.

  \item \textbf{Enum in \texttt{randomElement}.} Passing enum cases
    to \texttt{randomElement()} works because the factory's
    \texttt{definition()} return value is the raw PHP array that
    Eloquent receives before applying casts. Eloquent's
    \texttt{Gender} cast accepts either a \texttt{Gender} enum
    instance or its raw string value.
\end{itemize}

\subsection{DoctorScheduleFactory}

\begin{lstlisting}[
  style   = php,
  caption = {database/factories/DoctorScheduleFactory.php.},
  label   = {lst:factory-schedule},
]
<?php

namespace Database\Factories;

use App\Models\Clinic;
use App\Models\DoctorSchedule;
use App\Models\User;
use Illuminate\Database\Eloquent\Factories\Factory;

/** @extends Factory<DoctorSchedule> */
class DoctorScheduleFactory extends Factory
{
    protected $model = DoctorSchedule::class;

    public function definition(): array
    {
        return [
            'doctor_id'     => User::factory()->doctor(),
            'clinic_id'     => Clinic::factory(),
            'day_of_week'   => 1,         // Monday (1 = Monday in Carbon)
            'start_time'    => '08:00',
            'end_time'      => '17:00',
            'slot_duration' => 60,         // 60-minute slots
            'is_active'     => true,
        ];
    }

    public function inactive(): static
    {
        return $this->state(['is_active' => false]);
    }
}
\end{lstlisting}

\textbf{Fixed defaults for predictable tests.} Unlike most factories,
\phpclass{DoctorScheduleFactory} uses fixed rather than randomised
time values. This is deliberate: tests that exercise the slot
computation algorithm need to know exactly how many slots the schedule
will produce so they can write precise \texttt{assertCount()} assertions.

A schedule with \texttt{start\_time='08:00'}, \texttt{end\_time='17:00'},
and \texttt{slot\_duration=60} produces exactly
$\lfloor(17{:}00 - 08{:}00) / 60\text{ min}\rfloor = 9$ slots:
\texttt{08:00, 09:00, 10:00, 11:00, 12:00, 13:00, 14:00, 15:00, 16:00}.

\subsection{AppointmentFactory --- The Static Counter Technique}
\label{sec:appt-factory}

\phpclass{AppointmentFactory} is the most technically intricate factory
in the suite, requiring a solution to the composite unique index
constraint on \texttt{(doctor\_id, appointment\_date, start\_time)}.

\begin{lstlisting}[
  style   = php,
  caption = {database/factories/AppointmentFactory.php --- complete source.},
  label   = {lst:factory-appt},
]
<?php

namespace Database\Factories;

use App\Enums\AppointmentStatus;
use App\Models\Appointment;
use App\Models\Clinic;
use App\Models\Patient;
use App\Models\User;
use Illuminate\Database\Eloquent\Factories\Factory;

/** @extends Factory<Appointment> */
class AppointmentFactory extends Factory
{
    protected $model = Appointment::class;

    public function definition(): array
    {
        $clinic   = Clinic::factory()->create();
        $doctor   = User::factory()->doctor()->forClinic($clinic)->create();
        $patient  = Patient::factory()->create(['clinic_id' => $clinic->id]);
        $bookedBy = User::factory()->assistant()->forClinic($clinic)->create();

        // Static counter cycles start hours from 08 to 16 (9 distinct values)
        // to avoid the unique index on (doctor_id, appointment_date, start_time)
        // when creating multiple appointments for the same doctor in a batch.
        static $hour = 8;
        $startHour = ($hour++ % 9) + 8;   // 8, 9, 10, 11, 12, 13, 14, 15, 16
        $start     = sprintf('%02d:00', $startHour);
        $end       = sprintf('%02d:00', $startHour + 1);

        return [
            'clinic_id'        => $clinic->id,
            'doctor_id'        => $doctor->id,
            'patient_id'       => $patient->id,
            'booked_by'        => $bookedBy->id,
            'appointment_date' => now()->next('Monday')->format('Y-m-d'),
            'start_time'       => $start,
            'end_time'         => $end,
            'status'           => AppointmentStatus::PENDING,
            'notes'            => null,
        ];
    }

    public function confirmed(): static
    {
        return $this->state(['status' => AppointmentStatus::CONFIRMED]);
    }

    public function completed(): static
    {
        return $this->state(['status' => AppointmentStatus::COMPLETED]);
    }

    public function cancelled(): static
    {
        return $this->state(['status' => AppointmentStatus::CANCELLED]);
    }
}
\end{lstlisting}

\subsubsection{The Unique Index Problem}

The \texttt{appointments} table has a composite unique index:

\begin{lstlisting}[style=php, label={lst:unique-idx}]
$table->unique(['doctor_id', 'appointment_date', 'start_time']);
\end{lstlisting}

Without the static counter, if a test creates two appointment records
for the same doctor on the same date and both use the hardcoded
default \texttt{start\_time = '09:00'}, MySQL raises:

\begin{lstlisting}[
  style   = bash,
  caption = {Unique key violation error without the static counter.},
  label   = {lst:unique-error},
]
SQLSTATE[23000]: Integrity constraint violation: 1062
Duplicate entry 'uuid-2026-06-16-09:00:00' for key
'appointments_doctor_id_appointment_date_start_time_unique'
\end{lstlisting}

This error does not indicate a bug in the application code; it is a
test data construction problem.

\subsubsection{The Static Counter Solution}

The \texttt{static \$hour = 8} declaration inside \texttt{definition()}
creates a PHP class-static variable. In PHP, a \texttt{static} variable
inside a function retains its value between calls to that function within
the same process. It is different from a method-level local variable
(which is reset on every call) and from a class-level static property
(which is declared on the class, not inside a method).

The arithmetic \texttt{(\$hour++ \% 9) + 8} operates as follows:

\begin{table}[h!]
  \centering
  \caption{Static counter cycle: start hour produced on each factory call.}
  \label{tab:hour-cycle}
  \renewcommand{\arraystretch}{1.4}
  \begin{tabular}{@{}ccccc@{}}
    \toprule
    \rowcolor{PrimaryBlue}
    \tblhdr{Call} & \tblhdr{\$hour (before)} & \tblhdr{\$hour \% 9}
      & \tblhdr{+ 8} & \tblhdr{start\_time} \\
    \midrule
    1  & 8  & 8  & 16 & \texttt{16:00} \\
    \rowalt
    2  & 9  & 0  & 8  & \texttt{08:00} \\
    3  & 10 & 1  & 9  & \texttt{09:00} \\
    \rowalt
    4  & 11 & 2  & 10 & \texttt{10:00} \\
    5  & 12 & 3  & 11 & \texttt{11:00} \\
    \rowalt
    6  & 13 & 4  & 12 & \texttt{12:00} \\
    7  & 14 & 5  & 13 & \texttt{13:00} \\
    \rowalt
    8  & 15 & 6  & 14 & \texttt{14:00} \\
    9  & 16 & 7  & 15 & \texttt{15:00} \\
    \rowalt
    10 & 17 & 8  & 16 & \texttt{16:00} \\
    11 & 18 & 0  & 8  & \texttt{08:00} \\
    \bottomrule
  \end{tabular}
\end{table}

The cycle produces 9 distinct start times (08:00 through 16:00) and
then repeats. This covers any test case that creates up to 9 appointments
for the same doctor on the same date without collision. Tests that need
more than 9 appointments for the same doctor should provide explicit
\texttt{start\_time} overrides.

\begin{warningbox}[Static Counter Scope]
  The static counter is process-scoped, not test-scoped. If PHPUnit
  runs tests in the same PHP process (which is the default), the counter
  retains its value between tests. This is intentional: it means each
  test gets a fresh, never-before-used start time for any factory-created
  appointments. However, if the \phpclass{RefreshDatabase} trait rolls
  back a test's appointment records and the next test creates appointments
  for a \emph{different} doctor, the counter will start at whatever value
  it left off at, still producing valid unique times.
\end{warningbox}

\subsubsection{Self-Contained Relationship Creation in \texttt{definition()}}

Unlike other factories, \phpclass{AppointmentFactory::definition()} calls
\texttt{::create()} directly (not \texttt{::factory()}) for its related
models. This is because four related records (clinic, doctor, patient,
assistant) must be created \emph{before} the appointment's start time is
computed, and all four must be consistent with each other (same clinic).

The trade-off is that the factory is slower than a factory that uses
deferred factory references, and it always creates four extra records per
appointment in addition to the appointment itself. In tests that supply
all four related IDs as overrides, this behaviour is bypassed:

\begin{lstlisting}[style=php, caption={Overriding all relationships avoids the factory chain.}, label={lst:appt-override}]
// The definition() closure still runs, but the override array
// replaces clinic_id, doctor_id, patient_id, and booked_by
// AFTER the inner create() calls have already fired.
// To avoid the overhead, pass all four IDs explicitly:
$appt = Appointment::factory()->create([
    'clinic_id'        => $this->clinic->id,
    'doctor_id'        => $this->doctor->id,
    'patient_id'       => $this->patient->id,
    'booked_by'        => $this->assistant->id,
    'appointment_date' => now()->next('Monday')->format('Y-m-d'),
    'start_time'       => '08:00',
    'end_time'         => '09:00',
]);
// Note: the factory still creates the inner clinic/doctor/patient
// records because definition() runs first. To avoid this completely,
// use the constructor pattern if maximum speed is required.
\end{lstlisting}

\begin{tipbox}[Best Practice]
  In test \texttt{setUp()} methods, always create the fixture records
  once (clinic, doctor, assistant, patient) and reuse them across all
  tests in the class by passing their IDs as explicit overrides to the
  factory. This minimises database writes per test and makes the test
  data model explicit rather than hidden inside the factory.
\end{tipbox}

\subsection{PrescriptionFactory and PrescriptionItemFactory}

\begin{lstlisting}[
  style   = php,
  caption = {database/factories/PrescriptionFactory.php.},
  label   = {lst:factory-rx},
]
<?php

namespace Database\Factories;

use App\Models\Appointment;
use App\Models\Clinic;
use App\Models\Patient;
use App\Models\Prescription;
use App\Models\User;
use Illuminate\Database\Eloquent\Factories\Factory;

/** @extends Factory<Prescription> */
class PrescriptionFactory extends Factory
{
    protected $model = Prescription::class;

    public function definition(): array
    {
        return [
            'appointment_id' => Appointment::factory(),
            'doctor_id'      => User::factory()->doctor(),
            'patient_id'     => Patient::factory(),
            'clinic_id'      => Clinic::factory(),
            'diagnosis'      => fake()->sentence(),
            'notes'          => fake()->optional()->sentence(),
        ];
    }
}
\end{lstlisting}

\begin{lstlisting}[
  style   = php,
  caption = {database/factories/PrescriptionItemFactory.php.},
  label   = {lst:factory-rx-item},
]
<?php

namespace Database\Factories;

use App\Models\Prescription;
use App\Models\PrescriptionItem;
use Illuminate\Database\Eloquent\Factories\Factory;

/** @extends Factory<PrescriptionItem> */
class PrescriptionItemFactory extends Factory
{
    protected $model = PrescriptionItem::class;

    public function definition(): array
    {
        return [
            'prescription_id' => Prescription::factory(),
            'medicine_name'   => fake()->word()
                                 . ' '
                                 . fake()->numberBetween(5, 500)
                                 . 'mg',
            'dosage'      => fake()->randomElement([
                                '1 tablet', '2 tablets', '5ml', '10ml'
                             ]),
            'frequency'   => fake()->randomElement([
                                'Once daily', 'Twice daily', 'Three times daily'
                             ]),
            'duration'    => fake()->randomElement([
                                '3 days', '5 days', '7 days', '14 days'
                             ]),
            'notes'       => fake()->optional()->sentence(),
        ];
    }
}
\end{lstlisting}

\textbf{Notes on \texttt{fake()->optional()}.} The \texttt{optional()}
modifier wraps the subsequent faker call and returns \texttt{null} with
a configurable probability (default 50\%). This simulates the real-world
scenario where some prescription items include dispensing notes and
others do not, without writing a conditional in the factory definition.

%% ─────────────────────────────────────────────────────────────────────────
\section{Migration Execution and Schema Verification}
\label{sec:migration-execution}
%% ─────────────────────────────────────────────────────────────────────────

\begin{lstlisting}[
  style   = bash,
  caption = {Running and verifying migrations in development.},
  label   = {lst:migrate-cmds},
]
# Run all pending migrations in order
php artisan migrate

# Roll back the last batch of migrations (destructive)
php artisan migrate:rollback

# Show the current migration status
php artisan migrate:status

# Reset and re-run all migrations (destructive – clears all data)
php artisan migrate:fresh

# Reset, re-run migrations, and seed
php artisan migrate:fresh --seed

# Run only migrations for the test database (phpunit.xml overrides)
php artisan migrate --env=testing
\end{lstlisting}

\begin{warningbox}[Production Migrations]
  Never run \texttt{migrate:fresh} or \texttt{migrate:rollback} against
  a production database. These commands are irreversible and will destroy
  data. Use \texttt{migrate} (forward-only) in production and always take
  a database dump before any schema-altering migration run.
\end{warningbox}

The complete execution order of all domain migrations, from first to
last, is shown in Table~\ref{tab:migration-order}.

\begin{table}[h!]
  \centering
  \caption{Domain migration execution order.}
  \label{tab:migration-order}
  \renewcommand{\arraystretch}{1.4}
  \begin{tabularx}{\linewidth}{@{}C{0.5cm} L{7cm} X@{}}
    \toprule
    \rowcolor{PrimaryBlue}
    \tblhdr{\#} & \tblhdr{Migration File} & \tblhdr{Action} \\
    \midrule
    1 & \texttt{0001\_01\_01\_000000\_create\_users\_table}
        & Creates \texttt{users} (no FK yet), \texttt{password\_reset\_tokens},
          \texttt{sessions} \\
    \rowalt
    2 & \texttt{2026\_06\_05\_202852\_create\_clinics\_table}
        & Creates \texttt{clinics} \\
    3 & \texttt{2026\_06\_05\_202918\_create\_patients\_table}
        & Creates \texttt{patients} with \texttt{clinic\_id} FK \\
    \rowalt
    4 & \texttt{2026\_06\_05\_202954\_create\_doctor\_schedules\_table}
        & Creates \texttt{doctor\_schedules} \\
    5 & \texttt{2026\_06\_05\_203047\_create\_appointments\_table}
        & Creates \texttt{appointments} with composite unique index \\
    \rowalt
    6 & \texttt{2026\_06\_05\_203107\_create\_prescriptions\_table}
        & Creates \texttt{prescriptions} \\
    7 & \texttt{2026\_06\_05\_203128\_create\_prescription\_items\_table}
        & Creates \texttt{prescription\_items} \\
    \rowalt
    8 & \texttt{2026\_06\_11\_000001\_add\_foreign\_key\_to\_users\_clinic\_id}
        & Alters \texttt{users.clinic\_id} to UUID and adds FK \\
    9 & \texttt{2026\_06\_11\_000002\_add\_fulltext\_index\_to\_patients}
        & Adds FULLTEXT index on \texttt{patients(name, phone)} \\
    \bottomrule
  \end{tabularx}
\end{table}
